\begin{titlepage}
\centering
\vspace*{3.2cm}
{\huge\bfseries 数学 $\leftrightarrow$ 物理\par}
\vspace{0.6cm}
{\LARGE\bfseries 罗塞塔石碑：核心词典\par}
\vspace{1.2cm}
{\large The Rosetta Stone, Core Entries\par}
\vspace{0.4cm}
{\large 围绕 Kapustin--Witten 2006 与几何 Langlands 纲领\par}
\vspace{2.2cm}
{\large 核心词条 15 条 · 石碑总表 1 张 · 一页版石碑 1 张\par}
\vspace{0.3cm}
{\large 底本：arXiv:hep-th/0604151v3（231 页）\par}
\vspace{0.3cm}
{\large 对照综述：Witten 2008（arXiv:0802.0999）、Frenkel 2009（arXiv:0906.2747）\par}
\vfill
{\large 2026 年 11 月\par}
\vspace{0.8cm}
\end{titlepage}

\section*{前言：什么是「罗塞塔石碑」}
\addcontentsline{toc}{section}{前言：什么是「罗塞塔石碑」}

罗塞塔石碑是一块刻着同一段诏书的三种文字的石头：古埃及圣书体、世俗体、古希腊文。
它的价值不在任何一种文字本身，而在于**三种文字并排**——于是人们可以拿已知的一栏去破译未知的两栏。
安德烈·韦伊（André Weil）把这个比喻搬进数论，用来描述三栏彼此平行、逐句可译的图景：

\begin{center}\small
\begin{tabular}{lll}
\toprule
数论 & $\mathbb F_q$ 上的曲线 & 黎曼面 \\
\midrule
$\mathbb Q$ 与其有限扩张 & 函数域 $\mathbb F_q(C)$ & $\mathbb C$ 上的紧黎曼面 $C$ \\
素数 & 闭点 & 点 \\
Galois 表示 & $\ell$-进表示 & 局部系统（平坦丛） \\
自守表示 & 自守表示 & Hecke 特征层 / $\mathcal D$-模 \\
\bottomrule
\end{tabular}
\end{center}

韦伊自己的说法是：三栏我手上都只有残片，我知道三种语言各自的一些门道，
但栏与栏之间意义的差别常常出乎预料；这些年我只找到「词典的一小块一小块」。

本册要做的，是 KW06 这块「石碑」的**核心部件**：
Kapustin 与 Witten 在《电磁对偶与几何 Langlands 纲领》里给出的，是同一件事的两种语言——
左边是四维 $\mathcal N=4$ 超对称杨–米尔斯理论的物理，右边是黎曼面上的几何 Langlands 纲领的数学，
中间那条可逐句对译的带子，是**扭化的拓扑场论 + Hitchin 模空间 + 膜**。
本册不追求词条数量（那由详册《数学↔物理 罗塞塔词典》107 条承担），而是取其中
**十五个最核心的组件**，每条按百科词条的严格度写：精确定义、可核对的陈述、可手算的例子、出处与真伪状态。

\subsection*{怎么读}
\begin{enumerate}[label=\arabic*.,leftmargin=1.6em]
\item 每条固定七段：一句话 $\to$ 定义 $\to$ 定理/构造 $\to$ 物理侧翻译 $\to$ 可手算的例子 $\to$ 状态与出处 $\to$ 词典指针。
\item 「状态与出处」是区别本册与普通综述的地方：凡原文只做物理推导、数学上未证的，都标 argued 并说明缺哪一步。
\item 先读第 1 条（电磁对偶）与最后一条（几何 Langlands），中间的词条按需查。
\item 附录 A 是「石碑总表」：数学侧对象与物理侧对象的逐行对照，可当导航图用。
\end{enumerate}

\subsection*{与详册、与译本的关系}
\begin{itemize}[leftmargin=1.4em]
\item 详册《数学↔物理 罗塞塔词典》（107 条，7 轨道）负责「查」：一个词不懂，去那里看定义、例子、易错点。
\item 本册负责「立」：十五个核心组件各自的严格陈述与真伪状态。
\item 底本中译本《电磁对偶与几何 Langlands 纲领》中译注本（176 页，216 条译注）是引文与页码的依据：
  本册标注的「KW06 §x.y p.N」按中译本节号与原文印刷页给出。
\item 论证逻辑流程图见详册附录 B；本册附录 A 的总表与那张图互为经纬。
\end{itemize}

\subsection*{约定}
数学侧用蓝色小标题，物理侧用橙色；人名保留原文拼写；术语用详册统一译法
（电磁对偶、S-对偶、Langlands 对偶群 ${}^{L}G$、Hitchin 保留拉丁、特征膜 eigenbrane、
典范余各向同性膜、扭化 D-模）；公式编号只在需要引用时给。
